% =====================================================================
%  Paquetes. Cada uno con un comentario de para qué sirve.
% =====================================================================

% --- Idioma y codificación ---
\usepackage[T1]{fontenc}
\usepackage[utf8]{inputenc}
% es-tabla: "Tabla" en vez de "Cuadro"; es-noshorthands evita conflictos con TikZ
\usepackage[english, spanish, es-tabla, es-noshorthands, es-lcroman]{babel} % el último idioma es el principal
\usepackage{csquotes}           % comillas correctas con \enquote{...}

% --- Fuente (provisoria; cambiar según normativa) ---
\usepackage{lmodern}
% Alternativas tipo Times o Arial (descomentar UNA):
% \usepackage{newtxtext,newtxmath}   % parecida a Times New Roman
% \usepackage[scaled]{helvet}\renewcommand{\familydefault}{\sfdefault} % parecida a Arial

% --- Matemática ---
\usepackage{amsmath, amssymb, mathtools}
\usepackage{bm}                 % vectores/matrices en negrita: \bm{u}
\usepackage{icomma}             % coma decimal sin espacio extra en modo matemático

% --- Unidades (fundamental en ingeniería) ---
\usepackage{siunitx}
\sisetup{
  output-decimal-marker = {,},  % coma decimal (estilo español)
  group-separator = {.},
  group-minimum-digits = 4,
  per-mode = symbol,            % kN/m en vez de kN m^-1
  range-phrase = { a },
  list-final-separator = { y },
}

% --- Figuras, tablas y diagramas ---
\usepackage{graphicx}
\usepackage{booktabs}           % tablas profesionales (\toprule, \midrule, \bottomrule)
\usepackage{tabularx}
\usepackage{multirow}
\usepackage{longtable}          % tablas de varias páginas
\usepackage{caption}
\usepackage{subcaption}         % subfiguras (a), (b)
\usepackage{float}
\usepackage{tikz}               % diagramas dibujados con código
\usetikzlibrary{arrows.meta, positioning, calc, patterns, decorations.pathmorphing, shapes.geometric}
\usepackage{pgfplots}           % gráficos con código
\pgfplotsset{compat=1.18}

% --- Listas, colores, varios ---
\usepackage{enumitem}
\usepackage{xcolor}
\usepackage{setspace}           % interlineado
\usepackage{geometry}           % márgenes
\usepackage{fancyhdr}           % encabezados y pies
\usepackage{titlesec}           % formato de títulos
\usepackage{pdfpages}           % insertar PDFs (p.ej. hoja oficial firmada)
\usepackage{listings}           % código fuente en anexos

% --- Bibliografía ---
% Estilo provisorio APA. Para IEEE: style=ieee. Para numérico simple: style=numeric-comp
\usepackage[backend=biber, style=apa]{biblatex}

% --- Enlaces y referencias cruzadas (cargar al final) ---
\usepackage[hidelinks]{hyperref}
\usepackage[spanish, noabbrev]{cleveref}   % \cref{fig:x} -> "figura 2.1"
